% !TEX root = ../main.tex
	\section{Certified Safe Fallback}
	\label{section_safety_fallback}

    The hard constraint in \eqref{eq_important_task_constraint} requires every important task $\tau_i \in \boldsymbol{\tau}^{safe}$ to satisfy $Pr(r_i > D_i) \leq \Theta_i$ under the deployed configuration $\{\boldsymbol{P}, \boldsymbol{\mathcal{Q}}\}$.
    Because the optimizer of Sections~\ref{section_priority_opt}--\ref{section_config_opt} searches online over candidate priority assignments and QoS budgets, it must never \emph{ship} a configuration that violates this constraint---not even transiently, while a re-optimization is in progress or immediately after the environment $\textbf{E}_k$ has changed.
    The certified safe-fallback mechanism described below makes this important-task safety guarantee a built-in property of the framework rather than a conditional assumption on a user-supplied execution-time upper bound.

    \textbf{Certified fallback.}
    At every interval the scheduler runs a \emph{certified safe fallback}: the most recently deployed configuration $(\boldsymbol{P}^{\star},\boldsymbol{\mathcal{Q}}^{\star})$ that has been verified, through the probabilistic response-time analysis of \eqref{eq_prob_rta_in_opt}, to satisfy $Pr(r_i > D_i) \leq \Theta_i$ for every important task $\tau_i \in \boldsymbol{\tau}^{safe}$.
    Before the robot enters an unknown environment, an initial fallback is obtained by solving the feasibility problem posed by \eqref{eq_important_task_constraint} from a Deadline-Monotonic seed---the same important-first group lock used to seed the incremental search of Section~\ref{section_increment_pa}.
    If no feasible configuration exists for the current $\textbf{E}_k$, the system flags the environment as unschedulable and retains the previous fallback rather than deploying an unsafe one.

    \textbf{Adopt-or-retain.}
    Each candidate $\{\boldsymbol{P}^{(s)},\boldsymbol{\mathcal{Q}}^{(s)}\}$ produced by the incremental priority and QoS search (Sections~\ref{section_increment_pa} and~\ref{section_config_opt}) is scored against the current champion held in the RTA cache (Section~\ref{section_rta_cache}).
    A candidate is \emph{adopted}---i.e., promoted to the deployed fallback---only if it strictly improves the objective \eqref{eq_overall_obj} \emph{and} satisfies the important-task constraint $Pr(r_i > D_i) \leq \Theta_i$ for all $\tau_i \in \boldsymbol{\tau}^{safe}$, evaluated using the same response-time distribution $\rtDist{i}$ used for scoring.
    Candidates that violate the constraint, or that improve the objective at the cost of an important-task deadline-miss probability, are discarded, and the current fallback is retained.
    Optimization progress and safety are thereby decoupled: the optimizer may freely explore aggressive candidates, because any infeasible candidate is rejected in favor of the certified fallback.

    \begin{observation}[Important-task safety guarantee]
    \label{obs:safety_fallback}
    At every shipped re-optimization interval, the deployed configuration $(\boldsymbol{P}^{\star},\boldsymbol{\mathcal{Q}}^{\star})$ satisfies $Pr(r_i > D_i) \leq \Theta_i$ for every important task $\tau_i \in \boldsymbol{\tau}^{safe}$, regardless of whether the ongoing search has converged.
    \end{observation}

    The guarantee in Observation~\ref{obs:safety_fallback} holds by construction: only configurations that have passed the constraint check of \eqref{eq_important_task_constraint} are ever deployed, so an incomplete or interrupted optimization may degrade performance but can never violate the important-task safety bound.
    Crucially, this does \emph{not} require the designer to supply a conservative worst-case execution time (WCET) upper bound; the bound is certified directly against the probabilistic response-time analysis, so the guarantee is tight rather than pessimistic.
